\paragraph{\trCauchyRiemann{}算子。}\label{sub:cauchy-riemann-operators}
设 $S$ 是一个带边的\trRiemann\trSurface，$E\to S$ 是一个带有\trSymplecticStructure{}%
$\phi_E$ 与\trCompatible\trComplexStructure{}$J_E$ 的\trVectorBundle，而
$F\subset E|\partial S$ 是一个\trLagrange\trSubbundle。我们用 $C^\infty(S,E,F)$ 表示
那些限制到边界上后落于 $F$ 内的\trSection{}所成的空间，用
$C^\infty(S,\Omega^{0,1}_S\otimes E)$ 表示取值于 $E$ 的 $(0,1)$-形式所成的空间（不带
任何\trBoundaryCondition）。在\trSmooth{}度较低时，对一切 $2\le p<\infty$ 仍有%
\trWellDefined{}的空间 $W^{1,p}_{\mathrm{loc}}(S,E,F)$，即局部 $W^{1,p}$ 的\trSection{}所成
的空间（由于这样的\trSection{}是连续的，\trBoundaryCondition{}对 $p>2$ 显然有意义；
对 $p=2$，则经由限制映射 $W^{1,2}_{\mathrm{loc}}(S,E)\to
L^2_{\mathrm{loc}}(\partial S,E)$ 来定义它，例如见 [29]），以及局部 $L^p$ 的
$(0,1)$-形式所成的空间 $L^p_{\mathrm{loc}}(S,\Omega^{0,1}_S\otimes E)$。给定 $E$ 上
任意一个（不必是复的或辛的）\trConnection{}$\nabla$，人们把相配的%
\trCauchyRiemann{}算子定义为\trCovariantDifferentiation{}的 $(0,1)$ 部分：
\begin{equation}\label{eq:cr-operator}
\bar\partial_\nabla=\nabla^{0,1}\colon C^\infty(S,E,F)
\longrightarrow C^\infty(S,\Omega^{0,1}_S\otimes E).
\end{equation}
它可以扩张为 $W^{1,p}_{\mathrm{loc}}(S,E,F)\to
L^p_{\mathrm{loc}}(S,\Omega^{0,1}_S\otimes E)$。若 $S$ 是\trCompact{}的，该扩张便是一个%
\trFredholmOperator。
为了对非\trCompact{}的\trSurface{}上的\eqref{eq:cr-operator} 得到一套合理的分析理论，
需要假设 $S$ 具有\trStriplikeEnd，而且 $(E,F,\nabla)$ 在这些端上的\trAsymptotic{}行为
受到足够好的控制。精确的表述需要几个预备性的定义。对我们这个\trSurface{}$S$ 的每个
端 $\zeta$，配以一个\trVectorBundle{}$E_\zeta\to[0;1]$，它带有%
\trSymplecticStructure{}$\phi_{E_\zeta}$、与之\trCompatible{}的%
\trComplexStructure{}$J_{E_\zeta}$、一个\textbf{\trSymplectic\trConnection}（symplectic
connection）$\nabla_\zeta$，以及\trLagrange{}子空间
$\Lambda_{\zeta,k}\subset E_{\zeta,k}$（$k=0,1$）。我们称
$(E_\zeta,\phi_{E_\zeta},J_{E_\zeta},\nabla_\zeta,\Lambda_{\zeta,0},\Lambda_{\zeta,1})$
为 $\zeta$ 处的\textbf{\trLimitingDatum}（limiting datum）。把这个数据假设成%
\trNondegenerate{}的常会是方便的，其含义是：沿 $[0;1]$ 的
$\nabla_\zeta$-\trParallelTransport{}把 $\Lambda_{\zeta,0}$ 送到一个与
$\Lambda_{\zeta,1}$ \trTransverse{}的子空间。令 $\pi^\pm\colon Z^\pm\to[0;1]$ 为沿
$t$ 方向的\trProjection。我们要让丛 $E$ 带有特选的\trIsomorphism
\begin{equation}\label{eq:end-isomorphism}
\epsilon_\zeta^*E\cong(\pi^\pm)^*E_\zeta
\end{equation}
它们与诸\trLimitingDatum{}是\trAsymptotic{}地\trCompatible{}的。这意味着，如果我们认同
\eqref{eq:end-isomorphism} 的两边，则差 $\phi_E-\phi_{E_\zeta}$、$J_E-J_{E_\zeta}$、
$\nabla-\nabla_\zeta$ 随 $s\to\pm\infty$ 以指数速度衰减（比如说在 $C^1$
\trTopology{}下）；而且 $F_{\epsilon_\zeta(s,k)}\to\Lambda_{\zeta,k}$ 在同样的意义下%
\trConvergence（在 $(E_\zeta)_k$ 的\trLagrangianGrassmannian{}之内）。一个在诸端处
满足\eqref{eq:end-isomorphism} 的 $\bar\partial_\nabla$-算子称为%
\textbf{\trAdmissible}（admissible）的。另外，若一切\trLimitingDatum{}都是%
\trNondegenerate{}的，我们便说 $\bar\partial_\nabla$ 本身是\trNondegenerate{}的。
诸端上的\trTrivialization{}决定了一个沿 $s$ 方向为常值的\trMetric，我们用它来定义由
\emph{全局} $W^{1,p}$ 与 $L^p$ \trSection{}组成的空间 $W^{1,p}(S,E,F)$ 与
$L^p(S,\Omega^{0,1}_S\otimes E)$。于是\eqref{eq:cr-operator} 可扩张成一个连续算子
\begin{equation}\label{eq:global-extension}
\bar\partial_\nabla\colon W^{1,p}(S,E,F)
\longrightarrow L^p(S,\Omega^{0,1}_S\otimes E).
\end{equation}
若上文引入的\trNondegenerate{}条件成立，这些扩张都是\trFredholm{}的，其\trKernel{}与%
\trCokernel{}都与 $p$ 无关。更一般的做法是引入在诸端上带有某个权 $\epsilon\ge0$ 的
空间 $W^{1,p;\epsilon}$ 与 $L^{p;\epsilon}$（这些空间的元素是\trExponentialDecay{}的%
\trSection；当 $\epsilon=0$ 时它们就化归为前面考虑的寻常\trSobolevSpace）。对任何%
\trAdmissible{}的 $\bar\partial$-算子，存在一个 $\epsilon_0>0$，使得当 $\epsilon$ 取在
$(0;\epsilon_0)$ 中时，扩张
\begin{equation}\label{eq:cauchy-riemann-fredholm}
\bar\partial_\nabla\colon W^{1,p;\epsilon}(S,E,F)
\longrightarrow L^{p;\epsilon}(S,\Omega^{0,1}_S\otimes E)
\end{equation}
是\trFredholm{}的，且其\trKernel{}与\trCokernel{}与权无关（比较 [45, Proposition
4.1]）。若该\trCauchyRiemann{}算子是\trNondegenerate{}的，则同样的结论在
$\epsilon\in[0;\epsilon_0)$ 时也成立，这意味着在此情形，%
\eqref{eq:cauchy-riemann-fredholm} 的\trKernel{}与\trCokernel{}同
\eqref{eq:global-extension} 的一致。
有时我们会用到具有更高阶导数 $k>1$ 的\trSobolevSpace{}$W^{k,p}$，以及它们的带权版本
$W^{k,p;\epsilon}$，但同样的讨论对它们也成立。
